\documentclass[12pt,a4paper]{article}

\usepackage[utf8]{inputenc}
\usepackage[T1]{fontenc}
\usepackage[french]{babel}
\usepackage{amsmath,amssymb}
\usepackage{booktabs}
\usepackage{listings}
\usepackage[margin=2.5cm]{geometry}

\lstset{
  language=C++,
  basicstyle=\ttfamily\small,
  breaklines=true,
  frame=single,
  columns=fullflexible
}

\title{TP : parcours en largeur (BFS)\\d'un arbre binaire}
\author{Ali Arrai}
\date{28 octobre 2026}

\begin{document}
\maketitle

\section{Introduction}
Ce TP présente le parcours en largeur (BFS) appliqué à un arbre binaire,
programmé en C++. Contrairement aux parcours en profondeur, le BFS visite les
nœuds niveau par niveau : d'abord la racine, puis ses fils, puis les fils de
ses fils, et ainsi de suite. Il s'appuie sur une file (\texttt{queue}).

\section{Représentation de l'arbre}
Chaque nœud est un objet de la classe \texttt{noeud}, qui contient :
\begin{itemize}
  \item une valeur entière \texttt{data} ;
  \item un pointeur vers le fils gauche \texttt{fg} ;
  \item un pointeur vers le fils droit \texttt{fd}.
\end{itemize}
Des accesseurs \texttt{get\_fg}, \texttt{get\_fd} et des modificateurs
\texttt{set\_fg}, \texttt{set\_fd} permettent de relier les nœuds. La fonction
de parcours est déclarée \texttt{friend} afin de lire directement la valeur
d'un nœud.

\section{Algorithme}
Le principe est le suivant :
\begin{enumerate}
  \item si l'arbre est vide, on s'arrête ;
  \item on place la racine dans la file et on affiche sa valeur ;
  \item tant que la file n'est pas vide, on retire le nœud en tête ; si ce
        nœud a un fils gauche, on l'ajoute à la file et on affiche sa valeur ;
        on fait de même pour le fils droit.
\end{enumerate}
Les nœuds sont donc affichés dans l'ordre où ils entrent dans la file, c'est-à-dire
niveau par niveau et de gauche à droite.

\begin{lstlisting}
void parcourt_BFS(noeud* arbre) {
    if (arbre == nullptr) return;
    queue<noeud*> q;
    q.push(arbre);
    cout << arbre->data << ", ";
    while (!q.empty()) {
        noeud* current = q.front();
        q.pop();
        if (current->fg != nullptr) {
            q.push(current->fg);
            cout << current->fg->data << ", ";
        }
        if (current->fd != nullptr) {
            q.push(current->fd);
            cout << current->fd->data << ", ";
        }
    }
}
\end{lstlisting}

\section{Exemple et résultat}
Le programme principal construit un arbre de 12 nœuds de racine 7 :
\begin{itemize}
  \item 7 a pour fils 9 (gauche) et 67 (droit) ;
  \item 9 a pour fils 1 et 28 ; 28 a pour fils 29 et 2 ;
  \item 67 a pour fils 8 et 10 ; 8 a pour fils 57 et 40 ;
  \item 10 a un seul fils droit, 6.
\end{itemize}

L'exécution affiche l'ordre de parcours suivant :
\[
  7,\ 9,\ 67,\ 1,\ 28,\ 8,\ 10,\ 29,\ 2,\ 57,\ 40,\ 6.
\]

Cet ordre correspond bien aux niveaux de l'arbre :

\begin{table}[h!]
\centering
\begin{tabular}{cl}
\toprule
\textbf{Niveau} & \textbf{Nœuds visités} \\
\midrule
0 & 7 \\
1 & 9, 67 \\
2 & 1, 28, 8, 10 \\
3 & 29, 2, 57, 40, 6 \\
\bottomrule
\end{tabular}
\caption{Nœuds de l'arbre par niveau.}
\end{table}

\section{Analyse}
Chaque nœud entre une seule fois dans la file et en sort une seule fois, donc
le parcours coûte $O(n)$ pour un arbre de $n$ nœuds. La file contient au plus
un niveau de l'arbre, ce qui peut représenter jusqu'à $O(n)$ nœuds en mémoire
dans le pire cas, pour un arbre très large.

\section{Conclusion}
Ce TP montre comment une file permet de parcourir un arbre binaire en largeur,
en visitant les nœuds dans l'ordre de leur profondeur. Le résultat obtenu sur
l'exemple confirme ce comportement : les nœuds apparaissent niveau par niveau,
de gauche à droite.

\end{document}
